\documentclass[11pt,a4paper]{article}
\usepackage[T1]{fontenc}
\usepackage[utf8]{inputenc}
\usepackage{amsmath,amssymb,xcolor}
\usepackage[margin=25mm]{geometry}
\usepackage[colorlinks=true,urlcolor=blue]{hyperref}
\definecolor{nexorablue}{HTML}{2459B5}
\setlength{\parindent}{0pt}
\setlength{\parskip}{7pt}
\setlength{\emergencystretch}{2em}
\renewcommand{\familydefault}{\sfdefault}
\begin{document}

{\small\bfseries\color{nexorablue} NEXORA CPGE}\par
{\small MPSI / NEX-MES-03 / Corrigé}\par
{\LARGE\bfseries Mesurer une résistance par quotient}\par
\vspace{0.25em}\hrule\vspace{0.6em}

\section*{Estimation et propagation}
On trouve $R=250\,\mathrm\Omega$. Les contributions relatives valent $0{,}008$ et $0{,}010$. Leur combinaison vaut $\sqrt{0{,}008^2+0{,}010^2}\simeq0{,}0128$.

\section*{Bilan}
L’incertitude-type vaut environ $3{,}2\,\mathrm\Omega$. À deux chiffres significatifs, $R=(250{,}0\pm3{,}2)\,\mathrm\Omega$. L’intensité représente la plus grande part de la variance relative.


\end{document}
